% answers_short.tex —— 2020 年·北京卷 · 简短答案速查
% 由 tools/make_short_answers.py 自 answers.tex 的 \gksolhead/\gkanswer 提取，
% 并与 2020/2020_数学真题_北京_答案要点.md 逐题比对后生成，勿手改（改源 MD 后重跑）。

\begin{gknotice}
\textbf{使用说明}：本节只给最终答案便于快速核对，完整推导见解析版。答案为编写者独立推导，\textbf{非官方标准答案}。
\end{gknotice}

\gkpart{选择题}
\begin{center}\begin{tabular}{ccccccccccc}
\toprule
题号 & 1 & 2 & 3 & 4 & 5 & 6 & 7 & 8 & 9 & 10 \\
\midrule
答案 & D & B & C & D & A & D & B & B & C & A \\
\bottomrule
\end{tabular}\end{center}

\gkpart{填空题}
\begin{itemize}[leftmargin=2.6em]
  \item[\textbf{11.}] $(0,+\infty)$。需 $x+1\neq0$ 且 $x>0$，故定义域为 $(0,+\infty)$。
  \item[\textbf{12.}] $(3,0)$；$\sqrt3$。$c^2=6+3=9$，右焦点 $(3,0)$；渐近线 $y=\dfrac{x}{\sqrt2}$，距离 $\dfrac3{\sqrt3}=\sqrt3$。
  \item[\textbf{13.}] $\sqrt5$；$-1$。取 $A(0,0)$、$B(2,0)$、$C(2,2)$、$D(0,2)$，得 $P(2,1)$，故 $|\overrightarrow{PD}|=\sqrt{4+1}=\sqrt5$，$\overrightarrow{PB}\cdot\overrightarrow{PD}=(0,-1)\cdot(-2,1)=-1$。
  \item[\textbf{14.}] $\dfrac\pi2$（答案不唯一，$\varphi=2k\pi+\dfrac\pi2$ 均可）。振幅平方为 $\cos^2\varphi+(1+\sin\varphi)^2=2+2\sin\varphi=4$，故 $\sin\varphi=1$。
  \item[\textbf{15.}] ①②③。①区间平均下降速率甲大于乙；②$t_2$ 处甲曲线更陡且正由上方穿到下方；③$t_3$ 处两曲线均在达标线下方；④$[0,t_1]$ 段甲下降最缓，治理能力最弱，故④错。
\end{itemize}

\gkpart{解答题}
\begin{itemize}[leftmargin=2.6em]
  \item[\textbf{16.}] (1) 证明见解析。$\overrightarrow{BC_1}=\overrightarrow{AD_1}$，故 $BC_1\parallel AD_1\subset$ 平面 $AD_1E$，且 $BC_1$ 不在该平面内；(2) $\dfrac23$。平面 $AD_1E$ 的法向量取 $(1,2,-2)$，$\sin\theta=\dfrac{|\overrightarrow{AA_1}\cdot\boldsymbol n|}{|\overrightarrow{AA_1}||\boldsymbol n|}=\dfrac46=\dfrac23$
  \item[\textbf{17.}] (1) 条件①：$a=8$；条件②：$a=6$。条件①用余弦定理与 $a+b=11$ 得 $b=3$；条件②用正弦定理得 $a:b=\sin A:\sin B=6:5$；(2) 条件①：$\sin C=\dfrac{\sqrt3}2$，$S=6\sqrt3$；条件②：$\sin C=\dfrac{\sqrt7}4$，$S=\dfrac{15\sqrt7}4$
  \item[\textbf{18.}] (1) 男生 $\dfrac13$，女生 $\dfrac34$。由频率估计概率：$\dfrac{200}{600}=\dfrac13$，$\dfrac{300}{400}=\dfrac34$；(2) $\dfrac{13}{36}$。两名男生都支持而女生不支持的概率 $\left(\dfrac13\right)^2\cdot\dfrac14=\dfrac1{36}$，恰一名男生支持而女生支持的概率 $\mathrm C_2^1\cdot\dfrac13\cdot\dfrac23\cdot\dfrac34=\dfrac13$，相加即得；(3) $p_0>p_1$。一年级支持方案二的估计值为 $\dfrac{97}{192}>\dfrac12=p_0$，而 $p_0$ 是一年级与其他年级的加权平均，故其他年级的估计值 $p_1<p_0$
  \item[\textbf{19.}] (1) $y=-2x+13$。$f'(x)=-2x=-2$ 得切点 $(1,11)$，切线为 $y-11=-2(x-1)$；(2) 最小值为 $32$。切线 $y=-2tx+t^2+12$ 的两截距为 $\dfrac{t^2+12}{2t}$ 与 $t^2+12$，故 $S(t)=\dfrac{(t^2+12)^2}{4|t|}$，在 $|t|=2$ 处取最小值 $32$
  \item[\textbf{20.}] (1) $\dfrac{x^2}8+\dfrac{y^2}2=1$。由 $a=2b$ 得 $\dfrac4{4b^2}+\dfrac1{b^2}=1$，故 $b^2=2$，$a^2=8$；(2) $1$。设 $l:y=k(x+4)$，代入椭圆得 $(1+4k^2)x^2+32k^2x+64k^2-8=0$；算出 $y_P+y_Q=0$，即 $|PB|=|BQ|$
  \item[\textbf{21.}] (1) 不满足。取 $a_2=2$、$a_3=3$，则 $\dfrac{a_3^2}{a_2}=\dfrac92$ 不是正整数，不在 $\{a_n\}$ 中；(2) 同时满足性质①和性质②。$\dfrac{a_i^2}{a_j}=2^{2i-j-1}=a_{2i-j}$，且 $2i-j>0$；$a_n=2^{n-1}=\dfrac{a_{n-1}^2}{a_{n-2}}$（$n\geqslant3$）；(3) 证明见解析。先证各项同号；再用性质②定出 $(k,l)=(2,1)$，得 $a_3=a_2q$（$q=\dfrac{a_2}{a_1}$）；最后对 $n$ 归纳，由性质①给出上界 $a_{n+1}\leqslant a_1q^n$、性质②给出下界 $a_{n+1}\geqslant a_1q^n$，夹出 $a_{n+1}=a_nq$
\end{itemize}
